\chapter{以斯拉記}
\section{波斯王古列元年第一次歸回-建殿 1-6}
\begin{table}[H]
\centering 
\begin{tabular}{p{12cm}|p{1.5cm}}\hline\hline
\multicolumn{1}{c|}{古列詔書内容(1:1-4)}&执行过程\\\hline
神囑咐我在猶大的耶路撒冷為他建造殿宇&1:5-11\\\hline
在你們中間凡做他子民的，可以上猶大的耶路撒冷&2:1-67\\\hline
在耶路撒冷重建耶和華以色列神的殿&3-6 \\\hline
凡剩下的人要用金銀、財物、牲畜幫助，也要為耶路撒冷神的殿甘心獻禮物&2:68-70\\\hline
\end{tabular}
\end{table}

\subsection{居魯士(古列)王元年下詔猶大人返回建殿 1:1-4 }
波斯王居魯士元年，耶和華為要應驗藉耶利米口所說的話，就激動波斯王居魯士的心，使他下詔通告全國說： 2 「波斯王居魯士如此說：耶和華天上的神已將天下萬國賜給我，又囑咐我在猶大的耶路撒冷為他建造殿宇。 3 在你們中間凡做他子民的，可以上猶大的耶路撒冷，在耶路撒冷重建耶和華以色列神的殿，只有他是神。願神與這人同在！ 4 凡剩下的人，無論寄居何處，那地的人要用金銀、財物、牲畜幫助他，另外也要為耶路撒冷神的殿甘心獻上禮物。」

\subsection{波斯王古列盡返聖殿器物給猶大人 1:5-11 }
\begin{wraptable}{r}{5.5cm}
\centering 
\begin{tabular}{p{1cm}|p{1cm}|p{1cm}|p{1cm}}\hline\hline
器皿&數目& 器皿&數目\\\hline
金盤&30 &銀盤 &1000 \\\hline
刀&29 &金碗 &30 \\\hline
銀碗&410 &其它&1000 \\\hline
總數 &\multicolumn{3}{c}{5400} \\\hline
\end{tabular}
\end{wraptable}
於是，猶大和便雅憫的族長、祭司、利未人，就是一切被神激動他心的人，都起來，要上耶路撒冷去建造耶和華的殿。 6 他們四圍的人就拿銀器、金子、財物、牲畜、珍寶幫助他們(堅固他們的手)，另外還有甘心獻的禮物。 7 居魯士王也將耶和華殿的器皿拿出來，這器皿是尼布甲尼撒從耶路撒冷掠來，放在自己神之廟中的。 8 波斯王居魯士派庫官米提利達將這器皿拿出來，按數交給猶大的首領設巴薩。 9 器皿的數目記在下面：金盤三十個，銀盤一千個，刀二十九把， 10 金碗三十個，銀碗之次的四百一十個，別樣的器皿一千件， 11 金銀器皿共有五千四百件。被擄的人從巴比倫上耶路撒冷的時候，設巴薩將這一切都帶上來。


\subsection{被擄歸回的名單 2}
\subsubsection{首領的名單 2:1-2}
\begin{table}[H]
\centering 
\begin{tabular}{|p{2cm}|p{2cm}|p{2cm}|p{2cm}|p{2cm}|p{2cm}|}
\hline
\hline
所羅巴伯&耶書亞&尼希米  &西萊雅&利來雅&末底改\\\hline
必珊    &米斯拔&比革瓦伊&利宏  &巴拿  & \\\hline
\end{tabular}
%\caption{被擄歸回首領的名單}
\end{table}
巴比倫王尼布甲尼撒從前擄到巴比倫之猶大省的人，現在他們的子孫從被擄到之地回耶路撒冷和猶大，各歸本城。 2 他們是同著所羅巴伯、耶書亞、尼希米、西萊雅、利來雅、末底改、必珊、米斯拔、比革瓦伊、利宏、巴拿回來的。

\subsubsection{人民的名單 2:3-35}
\begin{table}[H]
\centering 
\begin{tabular}{|p{4cm}|p{1cm}|p{3cm}|p{1cm}|p{3cm}|p{1cm}|}\hline\hline
巴錄的子孫&2172 &示法提雅的子孫&372 &亞拉的子孫&775\\\hline
 巴哈摩押的後裔(耶書亞和約押的子孫)&2812&以攔的子孫  &1254&薩土的子孫&945\\\hline
薩改的子孫    &760&巴尼的子孫&642  &比拜的子孫  &623 \\\hline
押甲的子    &1222&亞多尼干的子孫&666  &比革瓦伊  &2056 \\\hline
亞丁的子孫    &454&亞特的後裔(希西家的子孫)&98  &比賽  &323 \\\hline
約拉的子孫    &112&哈順的子孫&223  &吉罷珥人  &95 \\\hline
伯利恆人    &123& 尼陀法人&56  &亞拿突人  &128 \\\hline
亞斯瑪弗人    &42&基列耶琳人、基非拉人、比錄人&743  &拉瑪人、迦巴人  &621 \\\hline
默瑪人人    &122& 伯特利人、艾人&223  &尼波人  &52 \\\hline
末必人    &156& 別的以攔子孫&1254  &哈琳的子孫  &320 \\\hline
羅德人、哈第人、阿挪人    &725&  耶利哥人&345  &西拿人  &3630 \\\hline
%伯利恆人    &123& 尼陀法人&56  &亞拿突人  &128 \\\hline
\end{tabular}
%\caption{被擄歸回人民的名單}
\end{table}
以色列人民的數目記在下面：巴錄的子孫二千一百七十二名， 4 示法提雅的子孫三百七十二名， 5 亞拉的子孫七百七十五名， 6 巴哈摩押的後裔，就是耶書亞和約押的子孫二千八百一十二名， 7 以攔的子孫一千二百五十四名， 8 薩土的子孫九百四十五名， 9 薩改的子孫七百六十名， 10 巴尼的子孫六百四十二名， 11 比拜的子孫六百二十三名， 12 押甲的子孫一千二百二十二名， 13 亞多尼干的子孫六百六十六名， 14 比革瓦伊的子孫二千零五十六名， 15 亞丁的子孫四百五十四名， 16 亞特的後裔，就是希西家的子孫九十八名， 17 比賽的子孫三百二十三名， 18 約拉的子孫一百一十二名， 19 哈順的子孫二百二十三名。 20 吉罷珥人九十五名， 21 伯利恆人一百二十三名， 22 尼陀法人五十六名， 23 亞拿突人一百二十八名， 24 亞斯瑪弗人四十二名， 25 基列耶琳人、基非拉人、比錄人共七百四十三名， 26 拉瑪人、迦巴人共六百二十一名， 27 默瑪人一百二十二名， 28 伯特利人、艾人共二百二十三名， 29 尼波人五十二名， 30 末必人一百五十六名， 31 別的以攔子孫一千二百五十四名， 32 哈琳的子孫三百二十名， 33 羅德人、哈第人、阿挪人共七百二十五名， 34 耶利哥人三百四十五名， 35 西拿人三千六百三十名。

\subsubsection{祭司的名單 2:36-39}
\begin{table}[H]
\centering 
\begin{tabular}{|p{4cm}|p{1cm}|p{4cm}|p{1cm}|}
\hline
\hline
耶書亞家耶大雅的子孫&973 &音麥的子孫&1052 \\\hline
巴施戶珥的子孫&1247&哈琳的子孫&1017\\\hline
\end{tabular}
%\caption{被擄歸回祭司的名單}
\end{table}
祭司，耶書亞家耶大雅的子孫九百七十三名， 37 音麥的子孫一千零五十二名， 38 巴施戶珥的子孫一千二百四十七名， 39 哈琳的子孫一千零一十七名。

\subsubsection{利未人的名單 2:40-58}
\begin{table}[H]
\centering 
\begin{tabular}{p{4cm}|p{10cm}|p{1cm}}
\hline
\hline
利未人&何達威雅的後裔，就是耶書亞和甲篾的子孫&74 \\\hline
歌唱的&亞薩的子孫&128\\\hline
守門的&沙龍的子孫、亞特的子孫、達們的子孫、亞谷的子孫、哈底大的子孫、朔拜的子孫&139\\\hline
尼提寧（就是殿役）& 西哈的子孫、哈蘇巴的子孫、答巴俄的子孫、基綠的子孫、西亞的子孫、巴頓的子孫、利巴拿的子孫、哈迦巴的子孫、亞谷的子孫、 哈甲的子孫、薩買的子孫、哈難的子孫、吉德的子孫、迦哈的子孫、利亞雅的子孫、利汛的子孫、尼哥大的子孫、迦散的子孫、烏撒的子孫、巴西亞的子孫、比賽的子孫、押拿的子孫、米烏寧的子孫、尼普心的子孫、巴卜的子孫、哈古巴的子孫、哈忽的子孫、巴洗律的子孫、米希大的子孫、哈沙的子孫、巴柯的子孫、西西拉的子孫、答瑪的子孫、尼細亞的子孫、哈提法的子孫。&\\\hline
所羅門僕人的後裔& 瑣太的子孫、瑣斐列的子孫、比路大的子孫、雅拉的子孫、達昆的子孫、吉德的子孫、示法提雅的子孫、哈替的子孫、玻黑列哈斯巴音的子孫、亞米的子孫&392\\\hline
\end{tabular}
%\caption{被擄歸回利未人的名單}
\end{table}
利未人，何達威雅的後裔，就是耶書亞和甲篾的子孫七十四名。 41 歌唱的，亞薩的子孫一百二十八名。 42 守門的，沙龍的子孫、亞特的子孫、達們的子孫、亞谷的子孫、哈底大的子孫、朔拜的子孫共一百三十九名。

\subsubsection{尼提寧人的名單 2:43-54}
尼提寧(殿役)，西哈的子孫、哈蘇巴的子孫、答巴俄的子孫、 44 基綠的子孫、西亞的子孫、巴頓的子孫、 45 利巴拿的子孫、哈迦巴的子孫、亞谷的子孫、 46 哈甲的子孫、薩買的子孫、哈難的子孫、 47 吉德的子孫、迦哈的子孫、利亞雅的子孫、 48 利汛的子孫、尼哥大的子孫、迦散的子孫、 49 烏撒的子孫、巴西亞的子孫、比賽的子孫、 50 押拿的子孫、米烏寧的子孫、尼普心的子孫、 51 巴卜的子孫、哈古巴的子孫、哈忽的子孫、 52 巴洗律的子孫、米希大的子孫、哈沙的子孫、 53 巴柯的子孫、西西拉的子孫、答瑪的子孫、 54 尼細亞的子孫、哈提法的子孫。

\subsubsection{所羅門僕人後裔的名單 2:55-58}
所羅門僕人的後裔，就是瑣太的子孫、瑣斐列的子孫、比路大的子孫、 56 雅拉的子孫、達昆的子孫、吉德的子孫、 57 示法提雅的子孫、哈替的子孫、玻黑列哈斯巴音的子孫、亞米的子孫。 58 尼提寧和所羅門僕人的後裔共三百九十二名。

%\clearpage
\subsubsection{其他不明身份的名單 2:59-63}
\begin{table}[H]
\centering 
\begin{tabular}{p{4cm}|p{10cm}|p{1cm}}
\hline
\hline
不能指明他們的宗族譜系是以色列人&第來雅的子孫、多比雅的子孫、尼哥大的子孫 & 652\\\hline
在族譜之中尋不著自己譜系的祭司&哈巴雅的子孫、哈哥斯的子孫、巴西萊的子孫；因為他們的先祖娶了基列人巴西萊的女兒為妻，所以起名叫巴西萊。&\\\hline
\end{tabular}
%\caption{不明身份的名單}
\end{table}
從特米拉、特哈薩、基綠、押但、音麥上來的，不能指明他們的宗族譜系是以色列人不是。 60 他們是第來雅的子孫、多比雅的子孫、尼哥大的子孫，共六百五十二名； 61 祭司中哈巴雅的子孫、哈哥斯的子孫、巴西萊的子孫（因為他們的先祖娶了基列人巴西萊的女兒為妻，所以起名叫巴西萊）， 62 這三家的人在族譜之中尋查自己的譜系，卻尋不著，因此算為不潔，不准供祭司的職任。 63 省長對他們說：「不可吃至聖的物，直到有用烏陵和土明決疑的祭司興起來。」

\subsubsection{總數 2:64-67}
\begin{table}[H]
\centering 
\begin{tabular}{p{1.5cm}|p{1.5cm}|p{1.5cm}|p{1.5cm}|p{2cm}|p{1.5cm}|p{1.5cm}|p{2cm}}
\hline
\hline
會眾 & 42360&僕婢& 7337&歌唱的男女&200&&  \\\hline
馬&736 &  騾子&245&駱駝&435&  驢&6720\\\hline
\end{tabular}
\caption{總數}
\end{table}
會眾共有四萬二千三百六十名。 65 此外還有他們的僕婢七千三百三十七名，又有歌唱的男女二百名。 66 他們有馬七百三十六匹，騾子二百四十五匹， 67 駱駝四百三十五隻，驢六千七百二十匹。

\subsection{奉獻 2:68-70}
\begin{table}[H]
\centering 
\begin{tabular}{p{2cm}|p{3cm}}
\hline
\hline
金子 & 61000達利克\\\hline
銀子&5000銀子\\\hline
祭司的禮服&100\\\hline
\end{tabular}
\caption{百姓甘心獻上禮物}
\end{table}
有些族長到了耶路撒冷耶和華殿的地方，便為神的殿甘心獻上禮物，要重新建造。 69 他們量力捐入工程庫的，金子六萬一千達利克，銀子五千彌拿，並祭司的禮服一百件。

70 於是，祭司、利未人、民中的一些人、歌唱的、守門的、尼提寧並以色列眾人，各住在自己的城裡。

\subsection{建造聖殿 3-6}
\subsubsection{築壇獻祭 3:1-7}	 
\subsubsubsection{築壇 3:1-2}
\begin{table}[H]
\centering 
\begin{tabular}{p{3cm}|p{3cm}|p{4cm}}
\hline
\hline
父親名字& 建壇的人名字& 其他人\\\hline
約薩達&耶書亞&和他的弟兄眾祭司\\\hline
撒拉鐵&所羅巴伯&與他的弟兄\\\hline
\end{tabular}
%\caption{參與建壇的人}
\end{table}
到了七月，以色列人住在各城，那時他們如同一人聚集在耶路撒冷。 2 約薩達的兒子耶書亞和他的弟兄眾祭司，並撒拉鐵的兒子所羅巴伯與他的弟兄，都起來建築以色列神的壇，要照神人摩西律法書上所寫的，在壇上獻燔祭。 3 

\subsubsubsection{獻祭 3:3-7}
\begin{table}[H]
\centering 
\begin{tabular}{|p{3cm}|p{6cm}|}
\hline
\hline
獻祭的名稱 & 獻祭的種類  \\\hline
燔祭       & 早晚獻      \\\hline
燔祭       &按數照例獻住棚節每日所當獻的\\\hline
燔祭 &月朔\\\hline
燔祭 &耶和華的一切聖節\\\hline
甘心祭 &各人向耶和華獻\\\hline
\end{tabular}
%\caption{在壇上所獻的祭}
\end{table}
他們在原有的根基上築壇，因懼怕鄰國的民；又在其上向耶和華早晚獻燔祭。 4 又照律法書上所寫的守住棚節，按數照例獻每日所當獻的燔祭。 5 其後獻常獻的燔祭，並在月朔與耶和華的一切聖節獻祭，又向耶和華獻各人的甘心祭。 6 從七月初一日起，他們就向耶和華獻燔祭，但耶和華殿的根基尚未立定。 7 他們又將銀子給石匠、木匠，把糧食、酒、油給西頓人、推羅人，使他們將香柏樹從黎巴嫩運到海裡，浮海運到約帕，是照波斯王居魯士所允准的。

%\subsection{重建聖殿 3:8-6:18}	 
\subsubsection{重立聖殿的根基 3:8-13}
\begin{table}[H]
\centering 
\begin{tabular}{p{3cm}|p{4cm}}
\hline
\hline
耶書亞&和他的子孫與弟兄  \\\hline
甲篾&和他的子孫\\\hline
利未人&希拿達的子孫與弟兄\\\hline
\end{tabular}
%\caption{建殿的督工}
\end{table}
百姓到了耶路撒冷神殿的地方第二年二月，撒拉鐵的兒子所羅巴伯、約薩達的兒子耶書亞，和其餘的弟兄，就是祭司、利未人，並一切被擄歸回耶路撒冷的人，都興工建造。又派利未人，從二十歲以外的，督理建造耶和華殿的工作。 9 於是猶大(2:40作「何達威雅」)的後裔，就是耶書亞和他的子孫與弟兄，甲篾和他的子孫，利未人希拿達的子孫與弟兄，都一同起來，督理那在神殿做工的人。 10 匠人立耶和華殿根基的時候，祭司皆穿禮服吹號，亞薩的子孫利未人敲鈸，照以色列王大衛所定的例都站著讚美耶和華。 11 他們彼此唱和，讚美稱謝耶和華說：「他本為善，他向以色列人永發慈愛！」他們讚美耶和華的時候，眾民大聲呼喊，因耶和華殿的根基已經立定。 12 然而有許多祭司、利未人、族長，就是見過舊殿的老年人，現在親眼看見立這殿的根基，便大聲哭號，也有許多人大聲歡呼， 13 甚至百姓不能分辨歡呼的聲音和哭號的聲音，因為眾人大聲呼喊，聲音聽到遠處。


\subsubsection{猶大和便雅憫的敵人參與建殿被拒 4:1-3}
猶大和便雅憫的敵人，聽說被擄歸回的人為耶和華以色列的神建造殿宇， 2 就去見所羅巴伯和以色列的族長，對他們說：「請容我們與你們一同建造，因為我們尋求你們的神，與你們一樣。自從亞述王以撒哈頓帶我們上這地以來，我們常祭祀神。」 3 但所羅巴伯、耶書亞和其餘以色列的族長對他們說：「我們建造神的殿與你們無干，我們自己為耶和華以色列的神協力建造，是照波斯王居魯士所吩咐的。」

\subsubsection{敵人破壞建殿,上本亞哈隨魯(亞達薛西)控告建殿 4:4-16}
那地的民就在猶大人建造的時候，使他們的手發軟，擾亂他們。 5 從波斯王居魯士年間直到波斯王大流士登基的時候，賄買謀士，要敗壞他們的謀算。 6 在亞哈隨魯才登基的時候，上本控告猶大和耶路撒冷的居民。亞達薛西年間，比施蘭、米特利達、他別和他們的同黨上本奏告波斯王亞達薛西，本章是用亞蘭文字亞蘭方言。 8 省長利宏、書記伸帥要控告耶路撒冷人，也上本奏告亞達薛西王： 9 「省長利宏、書記伸帥和同黨的底拿人、亞法薩提迦人、他毗拉人、亞法撒人、亞基衛人、巴比倫人、書珊迦人、底亥人、以攔人， 10 和尊大的亞斯那巴所遷移，安置在撒馬利亞城並大河西一帶地方的人等， 11 上奏亞達薛西王說：河西的臣民云云。 12 王該知道，從王那裡上到我們這裡的猶大人已經到耶路撒冷，重建這反叛惡劣的城，築立根基，建造城牆。 13 如今王該知道，他們若建造這城，城牆完畢就不再與王進貢、交課、納稅，終久王必受虧損。 14 我們既食御鹽，不忍見王吃虧，因此奏告於王。 15 請王考察先王的實錄，必在其上查知這城是反叛的城，於列王和各省有害，自古以來其中常有悖逆的事，因此這城曾被拆毀。 16 我們謹奏王知，這城若再建造，城牆完畢，河西之地王就無份了。」


\subsubsection{亞達薛西王王令停建神殿 4:17-24}
\begin{table}[!hbp]
\centering 
\begin{tabular}{|p{1.5cm}|p{3cm}|p{8cm}|}
\hline
\hline
第一次 & 亞哈隨魯才登基&那地的民上本控告猶大和耶路撒冷的居民  \\\hline
第二次 & 亞達薛西年間  &比施蘭、米特利達、他別，和他們的同黨上本奏告\\\hline
第三次 & 亞達薛西      &省長利宏、書記伸帥也上本奏告   \\\hline
\end{tabular}
\caption{參與上本奏告的人}
\end{table}
那時王諭復省長利宏、書記伸帥，和他們的同黨，就是住撒馬利亞並河西一帶地方的人，說：「願你們平安云云。 18 你們所上的本，已經明讀在我面前。 19 我已命人考查，得知此城古來果然背叛列王，其中常有反叛悖逆的事。 20 從前耶路撒冷也有大君王統管河西全地，人就給他們進貢、交課、納稅。 21 現在你們要出告示命這些人停工，使這城不得建造，等我降旨。 22 你們當謹慎，不可遲延，為何容害加重使王受虧損呢？」

23 亞達薛西王的上諭讀在利宏和書記伸帥並他們的同黨面前，他們就急忙往耶路撒冷去見猶大人，用勢力強迫他們停工。 24 於是，在耶路撒冷神殿的工程就停止了，直停到波斯王大流士第二年。


\begin{table}[H]
\centering 
\begin{tabular}{p{2cm}|p{2.5cm}|p{2cm}|p{2cm}|p{2cm}}
\hline
\hline
底拿人&亞法薩提迦人& 他毗拉人&亞法撒人&亞基衛人 \\\hline
巴比倫人&書珊迦人&底亥人&以攔人&\\\hline
\multicolumn{5}{|c|}{亞斯那巴所遷移、安置在撒瑪利亞城，並大河西一帶地方的人}\\\hline
\end{tabular}
\caption{省長利宏、書記伸帥的同黨}
\end{table}


\subsubsection{先知哈該,撒迦利亞幫助,建殿工作恢復 5:1-5}
\begin{table}[!hbp]
\centering 
\begin{tabular}{|p{4cm}|p{5cm}|}
\hline
\hline
先知& 領袖 \\\hline
先知哈該&撒拉鐵的兒子所羅巴伯\\\hline
易多的孫子撒迦利亞&約薩達的兒子耶書亞\\\hline
\end{tabular}
\caption{建殿工作恢復}
\end{table}
那時，先知哈該和易多的孫子撒迦利亞奉以色列神的名，向猶大和耶路撒冷的猶大人說勸勉的話。 2 於是，撒拉鐵的兒子所羅巴伯和約薩達的兒子耶書亞都起來，動手建造耶路撒冷神的殿，有神的先知在那裡幫助他們。 3 當時河西的總督達乃和示他波斯乃並他們的同黨來問說：「誰降旨讓你們建造這殿，修成這牆呢？」 4 我們便告訴他們建造這殿的人叫什麼名字。 5 神的眼目看顧猶大的長老，以致總督等沒有叫他們停工，直到這事奏告大流士得著他的回諭。

\subsubsection{河西的總督達乃並同黨上本大利烏王 5:6-17}
\begin{table}[!hbp]
\centering 
\begin{tabular}{|p{1.5cm}|p{3cm}|p{8cm}|}
\hline
\hline
第四次 & 大利烏王 (大流士）&河西的總督達乃和示他波斯乃，並他們的同黨，就是住河西的亞法薩迦人  \\\hline
\end{tabular}
\caption{參與上本奏告的人}
\end{table}
河西的總督達乃和示他波斯乃，並他們的同黨，就是住河西的亞法薩迦人，上本奏告大流士王。 7 本上寫著說：「願大流士王諸事平安！ 8 王該知道，我們往猶大省去，到了至大神的殿，這殿是用大石建造的，梁木插入牆內，工作甚速，他們手下亨通。 9 我們就問那些長老說：『誰降旨讓你們建造這殿，修成這牆呢？』 10 又問他們的名字，要記錄他們首領的名字，奏告於王。 11 他們回答說：『我們是天地之神的僕人，重建前多年所建造的殿，就是以色列的一位大君王建造修成的。 12 只因我們列祖惹天上的神發怒，神把他們交在迦勒底人巴比倫王尼布甲尼撒的手中，他就拆毀這殿，又將百姓擄到巴比倫。 13 然而巴比倫王居魯士元年，他降旨允准建造神的這殿。 14 神殿中的金銀器皿，就是尼布甲尼撒從耶路撒冷的殿中掠去帶到巴比倫廟裡的，居魯士王從巴比倫廟裡取出來，交給派為省長的，名叫設巴薩， 15 對他說：「可以將這些器皿帶去，放在耶路撒冷的殿中，在原處建造神的殿。」 16 於是，這設巴薩來建立耶路撒冷神殿的根基。這殿從那時直到如今，尚未造成。』

「現在王若以為美，請察巴比倫王的府庫，看居魯士王降旨允准在耶路撒冷建造神的殿沒有。王的心意如何，請降旨曉諭我們。」


\subsubsection{大利烏（大流士）王得先王之詔 6:1-5}
於是大流士王降旨，要尋察典籍庫內，就是在巴比倫藏寶物之處。 2 在瑪代省亞馬他城的宮內尋得一卷，其中記著說， 3 居魯士王元年，他降旨論到耶路撒冷神的殿：要建造這殿為獻祭之處，堅立殿的根基，殿高六十肘，寬六十肘， 4 用三層大石頭，一層新木頭，經費要出於王庫。 5 並且神殿的金銀器皿，就是尼布甲尼撒從耶路撒冷的殿中掠到巴比倫的，要歸還帶到耶路撒冷的殿中，各按原處放在神的殿裡。

\subsubsection{大利烏（大流士）王命勿阻建殿 6:6-12}
現在河西的總督達乃和示他波斯乃，並你們的同黨，就是住河西的亞法薩迦人，你們當遠離他們。 7 不要攔阻神殿的工作，任憑猶大人的省長和猶大人的長老在原處建造神的這殿。 8 我又降旨，吩咐你們向猶大人的長老為建造神的殿當怎樣行，就是從河西的款項中，急速撥取貢銀做他們的經費，免得耽誤工作。 9 他們與天上的神獻燔祭所需用的公牛犢、公綿羊、綿羊羔，並所用的麥子、鹽、酒、油，都要照耶路撒冷祭司的話，每日供給他們，不得有誤， 10 好叫他們獻馨香的祭給天上的神，又為王和王眾子的壽命祈禱。 11 我再降旨，無論誰更改這命令，必從他房屋中拆出一根梁來，把他舉起，懸在其上，又使他的房屋成為糞堆。 12 若有王和民伸手更改這命令，拆毀這殿，願那使耶路撒冷的殿作為他名居所的神將他們滅絕。我大流士降這旨意，當速速遵行。

\subsubsection{大利烏（大流士）第六年建殿完成 6:13-15}
於是，河西總督達乃和示他波斯乃並他們的同黨，因大流士王所發的命令，就急速遵行。 14 猶大長老因先知哈該和易多的孫子撒迦利亞所說勸勉的話，就建造這殿，凡事亨通。他們遵著以色列神的命令和波斯王居魯士、大流士、亞達薛西的旨意，建造完畢。 15 大流士王第六年，亞達月初三日，這殿修成了。

\subsubsection{行獻殿禮 6:16-18}
\begin{table}[!hbp]
\centering 
\begin{tabular}{p{1.5cm}|p{1.5cm}|p{1.5cm}|p{1.5cm}}
\hline
\hline
公牛 & 100 &公綿羊& 200 \\\hline
綿羊羔& 400&公山羊 &12  \\\hline
\end{tabular}
\caption{獻殿禮的祭物}
\end{table}
以色列的祭司和利未人，並其餘被擄歸回的人，都歡歡喜喜地行奉獻神殿的禮。 17 行奉獻神殿的禮就獻公牛一百隻，公綿羊二百隻，綿羊羔四百隻；又照以色列支派的數目獻公山羊十二隻，為以色列眾人做贖罪祭。 18 且派祭司和利未人按著班次在耶路撒冷侍奉神，是照摩西律法書上所寫的。

\subsubsection{正月十四日百姓守逾越節 6:19-22}	 
正月十四日，被擄歸回的人守逾越節。 20 原來，祭司和利未人一同自潔，無一人不潔淨。利未人為被擄歸回的眾人和他們的弟兄眾祭司，並為自己宰逾越節的羊羔。 21 從擄到之地歸回的以色列人，和一切除掉所染外邦人汙穢歸附他們，要尋求耶和華以色列神的人，都吃這羊羔。 22 歡歡喜喜地守除酵節七日，因為耶和華使他們歡喜，又使亞述王的心轉向他們，堅固他們的手，做以色列神殿的工程。
	 
\section{亞達薛西年間第二次歸回-教訓律法 7-8}
\subsection{介绍以斯拉 7:1-10}
\subsubsection{以斯拉的家譜 7:1-5}
\begin{table}[H]
\centering 
\begin{tabular}{p{1.5cm}|p{1.5cm}|p{1.2cm}|p{1.2cm}|p{1.2cm}|p{1.2cm}|p{1.5cm}|p{1.2cm}|p{1.2cm}}\hline\hline
亞倫 & 以利亞撒 &非尼哈& 亞比書& 布基 & 烏西& 西拉希雅 &米拉約&\\\hline
亞撒利雅& 亞瑪利雅&亞希突 &撒督 & 沙龍 & 希勒家 &亞撒利雅& 西萊雅 &以斯拉\\\hline
\end{tabular}
%\caption{以斯拉的家譜}
\end{table}
這事以後，波斯王亞達薛西年間，有個以斯拉，他是西萊雅的兒子。西萊雅是亞撒利雅的兒子，亞撒利雅是希勒家的兒子， 2 希勒家是沙龍的兒子，沙龍是撒督的兒子，撒督是亞希突的兒子， 3 亞希突是亞瑪利雅的兒子，亞瑪利雅是亞撒利雅的兒子，亞撒利雅是米拉約的兒子， 4 米拉約是西拉希雅的兒子，西拉希雅是烏西的兒子，烏西是布基的兒子， 5 布基是亞比書的兒子，亞比書是非尼哈的兒子，非尼哈是以利亞撒的兒子，以利亞撒是大祭司亞倫的兒子。

\subsubsection{以斯拉率民第二次歸回,將律例、典章教訓以色列人 7:6-10}
\begin{wraptable}{r}{6.5cm}
\centering 
\begin{tabular}{p{3.5cm}|p{2.5cm}}\hline\hline
王第七年正月初一日&從巴比倫起程  \\\hline
王第七年五月初一日 &到了耶路撒冷 \\\hline
\end{tabular}
%\caption{第二次歸回}
\end{wraptable}
這以斯拉從巴比倫上來，他是敏捷的文士，通達耶和華以色列神所賜摩西的律法書。王允准他一切所求的，是因耶和華他神的手幫助他。 7 亞達薛西王第七年，以色列人、祭司、利未人、歌唱的、守門的、尼提寧，有上耶路撒冷的。 8 王第七年五月，以斯拉到了耶路撒冷。 9 正月初一日他從巴比倫起程，因他神施恩的手幫助他，五月初一日就到了耶路撒冷。以斯拉定志考究遵行耶和華的律法，又將律例、典章教訓以色列人。

\subsection{亞達薛西王的諭旨和以斯拉的赞美 7:11-28}
\subsubsection{賜給以斯拉的諭旨 7:11-20}
\begin{table}[H]
\centering 
\begin{tabular}{|p{3cm}|p{8cm}|}\hline\hline
\multicolumn{2}{|l|}{國中的以色列人、祭司、利未人，凡甘心上耶路撒冷去的旨准他們與你同去}\\\hline
\multicolumn{2}{|l|}{所交給你神殿中使用的器皿，你要交在耶路撒冷神面前}\\\hline
\multicolumn{2}{|l|}{照你手中神的律法書察問猶大和耶路撒冷的景況} \\\hline
\multirow{4}{6em}{攜帶金銀來源} &王和謀士甘心獻給住耶路撒冷、以色列神\\
                          &你在巴比倫全省所得的金銀\\
                          &百姓、祭司樂意獻給耶路撒冷他們神殿的禮物\\
                          &神殿裡若再有需用的經費，可以從王的府庫裡支取\\\hline
\multirow{2}{8em}{金銀用途} &買公牛,公綿羊,綿羊羔,素祭奠祭，獻在神殿壇上\\ 
&剩下的金銀，你和你的弟兄看著怎樣好，就怎樣用\\\hline
\end{tabular}
%\caption{王賜給以斯拉的諭旨}
\end{table}
祭司以斯拉是通達耶和華誡命和賜以色列之律例的文士。亞達薛西王賜給他諭旨，上面寫著說： 12 「諸王之王亞達薛西，達於祭司以斯拉，通達天上神律法大德的文士，云云。 13 住在我國中的以色列人、祭司、利未人，凡甘心上耶路撒冷去的，我降旨准他們與你同去。
 14 王與七個謀士既然差你去，照你手中神的律法書察問猶大和耶路撒冷的景況， 

15 又帶金銀，就是王和謀士甘心獻給住耶路撒冷以色列神的， 16 並帶你在巴比倫全省所得的金銀和百姓、祭司樂意獻給耶路撒冷他們神殿的禮物， 17 所以你當用這金銀急速買公牛、公綿羊、綿羊羔和同獻的素祭奠祭之物，獻在耶路撒冷你們神殿的壇上。 18 剩下的金銀，你和你的弟兄看著怎樣好，就怎樣用，總要遵著你們神的旨意。 19 所交給你神殿中使用的器皿，你要交在耶路撒冷神面前。 20 你神殿裡若再有需用的經費，你可以從王的府庫裡支取。

\subsubsection{賜給河西庫官的諭旨 7:21-26}
\begin{table}[!hbp]
\centering 
\begin{tabular}{|p{6cm}|p{5cm}|}
\hline
\hline
銀子& 直到一百他連得\\
麥子&一百柯珥\\
酒&一百罷特\\
油&一百罷特\\
鹽&不計其數\\\hline
祭司、利未人、歌唱的、守門的，和尼提寧，並在神殿當差的人& 不可叫他們進貢，交課，納稅\\\hline
\multicolumn{2}{|l|}{將所有明白你神律法的人立為士師、審判官，治理河西的百姓}\\\hline
\multicolumn{2}{|l|}{不遵行你神律法和王命令的人就當速速定他的罪，或治死，或充軍，或抄家，或囚禁}\\\hline
\end{tabular}
\caption{王賜給河西庫官的諭旨}
\end{table}
我，亞達薛西王，又降旨於河西的一切庫官說：通達天上神律法的文士祭司以斯拉無論向你們要什麼，你們要速速地備辦， 22 就是銀子直到一百他連得、麥子一百柯珥、酒一百罷特、油一百罷特、鹽不計其數，也要給他。 

23 凡天上之神所吩咐的，當為天上神的殿詳細辦理，為何使憤怒臨到王和王眾子的國呢？ 

24 我又曉諭你們：至於祭司、利未人、歌唱的、守門的和尼提寧並在神殿當差的人，不可叫他們進貢、交課、納稅。

 25 以斯拉啊，要照著你神賜你的智慧，將所有明白你神律法的人立為士師、審判官，治理河西的百姓，使他們教訓一切不明白神律法的人。 26 凡不遵行你神律法和王命令的人，就當速速定他的罪，或治死，或充軍，或抄家，或囚禁。」

\subsubsection{以斯拉稱頌耶和華 7:27-28}
以斯拉說：「耶和華我們列祖的神是應當稱頌的！因他使王起這心意修飾耶路撒冷耶和華的殿， 28 又在王和謀士並大能的軍長面前施恩於我。因耶和華我神的手幫助我，我就得以堅強，從以色列中招聚首領與我一同上來。」

\subsection{歸回人員名單 8:1-20}
\subsubsection{記偕以斯拉返自巴比倫之人 8:1-14}
\begin{table}[H]
\centering 
\begin{tabular}{p{5cm}|p{4.5cm}|p{2cm}}\hline\hline
%父親名字& 建壇的人名字& 其他人\\\hline
非尼哈的子孫&革順& \\\hline
以他瑪的子孫&但以理& \\\hline
大衛的子孫& & \\\hline
巴錄的後裔(示迦尼的子孫) &撒迦利亞  &男丁150\\\hline
巴哈摩押的子孫& 西拉希雅的兒子以利約乃  &男丁200\\\hline
示迦尼的子孫& 雅哈悉的兒子  &男丁300\\\hline
亞丁的子孫& 約拿單的兒子以別  &男丁50\\\hline
以攔的子孫& 亞他利雅的兒子耶篩亞  &男丁70\\\hline
示法提雅的子孫& 米迦勒的兒子西巴第雅  &男丁80\\\hline
約押的子孫& 耶歇的兒子俄巴底亞  &男丁210\\\hline
示羅密的子孫& 約細斐的兒子  &男丁160\\\hline
比拜的子孫& 比拜的兒子撒迦利亞 &男丁28\\\hline
押甲的子孫&哈加坦的兒子約哈難，&男丁110\\\hline
亞多尼干的子孫，就是末尾的&以利法列、耶利、示瑪雅&男丁60\\\hline
比革瓦伊的子孫&烏太和撒布&男丁70\\\hline
\end{tabular}
%\caption{歸回人員名單}
\end{table}
當亞達薛西王年間，同我從巴比倫上來的人，他們的族長和他們的家譜記在下面： 2 屬非尼哈的子孫有革順；屬以他瑪的子孫有但以理；屬大衛的子孫有哈突； 3 屬巴錄的後裔，就是示迦尼的子孫，有撒迦利亞，同著他按家譜計算，男丁一百五十人； 4 屬巴哈摩押的子孫有西拉希雅的兒子以利約乃，同著他有男丁二百； 5 屬示迦尼的子孫有雅哈悉的兒子，同著他有男丁三百； 6 屬亞丁的子孫有約拿單的兒子以別，同著他有男丁五十； 7 屬以攔的子孫有亞他利雅的兒子耶篩亞，同著他有男丁七十； 8 屬示法提雅的子孫有米迦勒的兒子西巴第雅，同著他有男丁八十； 9 屬約押的子孫有耶歇的兒子俄巴底亞，同著他有男丁二百一十八； 10 屬示羅密的子孫有約細斐的兒子，同著他有男丁一百六十； 11 屬比拜的子孫有比拜的兒子撒迦利亞，同著他有男丁二十八； 12 屬押甲的子孫有哈加坦的兒子約哈難，同著他有男丁一百一十； 13 屬亞多尼干的子孫，就是末尾的，他們的名字是以利法列、耶利、示瑪雅，同著他們有男丁六十； 14 屬比革瓦伊的子孫有烏太和撒布，同著他們有男丁七十。

\subsubsection{召利未人與尼提寧人至 8:16-20}
\begin{table}[H]
\centering 
\begin{tabular}{p{1.5cm}|p{13.5cm}}
\hline
\hline
%父親名字& 建壇的人名字& 其他人\\\hline
利未首領&以利以謝、亞列、示瑪雅、以利拿單、雅立、以利拿單、拿單、撒迦利亞、米書蘭 \\\hline
教習&約雅立和以利拿單 \\\hline
\end{tabular}
%\caption{亞哈瓦河邊召利未人}
\end{table}
我招聚這些人在流入亞哈瓦的河邊，我們在那裡住了三日。我查看百姓和祭司，見沒有利未人在那裡， 16 就召首領以利以謝、亞列、示瑪雅、以利拿單、雅立、以利拿單、拿單、撒迦利亞、米書蘭，又召教習約雅立和以利拿單。 17 我打發他們往迦西斐雅地方去見那裡的首領易多，又告訴他們當向易多和他的弟兄尼提寧說什麼話，叫他們為我們神的殿帶使用的人來。 18 蒙我們神施恩的手幫助我們，他們在以色列的曾孫、利未的孫子、抹利的後裔中帶一個通達人來，還有示利比和他的眾子與弟兄共一十八人。 19 又有哈沙比雅，同著他有米拉利的子孫耶篩亞並他的眾子和弟兄，共二十人。 20 從前大衛和眾首領派尼提寧服侍利未人，現在從這尼提寧中也帶了二百二十人來，都是按名指定的。

\begin{table}[H]
\centering 
\begin{tabular}{p{10.5cm}|p{1.5cm}}
\hline
\hline
%父親名字& 建壇的人名字& 其他人\\\hline
以色列的曾孫、利未的孫子、抹利的後裔中&1通達人\\\hline
示利比和他的眾子與弟兄&118 \\\hline
哈沙比雅，同著他有米拉利的子孫耶篩亞，並他的眾子和弟兄&20 \\\hline
從前大衛和眾首領派尼提寧服事利未人&220 \\\hline
\end{tabular}
\caption{在迦西斐雅召利未人}
\end{table}


\subsection{第二批百姓回归 8:21-23}
\subsubsection{以斯拉領百姓在亞哈瓦河邊禁食禱告 8:21-23}
那時我在亞哈瓦河邊宣告禁食，為要在我們神面前克苦己心，求他使我們和婦人孩子並一切所有的都得平坦的道路。 
22 我求王撥步兵馬兵幫助我們抵擋路上的仇敵本以為羞恥，因我曾對王說：「我們神施恩的手必幫助一切尋求他的，但他的能力和憤怒必攻擊一切離棄他的。」 
23 所以我們禁食祈求我們的神，他就應允了我們。


\subsubsection{祭司權衡金銀器物，攜帶的聖殿禮物 8:24-30}
\begin{table}[H]
\centering 
\begin{tabular}{|p{4.5cm}|p{3.5cm}|}
\hline
\hline
銀子&650他連得\\\hline
銀器&100他連得 \\\hline
金子&100他連得\\\hline
金碗&20個, 重1000達利克\\\hline
上等光銅的器皿,寶貴如金&2個\\\hline
\end{tabular}
\caption{歸回時攜帶的聖殿禮物}
\end{table}
我分派祭司長十二人，就是示利比、哈沙比雅和他們的弟兄十人， 25 將王和謀士、軍長並在那裡的以色列眾人為我們神殿所獻的金銀和器皿，都稱了交給他們。 26 我稱了交在他們手中的，銀子有六百五十他連得，銀器重一百他連得，金子一百他連得， 27 金碗二十個，重一千達利克，上等光銅的器皿兩個，寶貴如金。 28 我對他們說：「你們歸耶和華為聖，器皿也為聖，金銀是甘心獻給耶和華你們列祖之神的。 29 你們當警醒看守，直到你們在耶路撒冷耶和華殿的庫內，在祭司長和利未族長並以色列的各族長面前過了秤。」 30 於是，祭司、利未人按著分量接受金銀和器皿，要帶到耶路撒冷我們神的殿裡。

\begin{table}[H]
\centering 
\begin{tabular}{|p{2cm}|p{5cm}|}
\hline
\hline
\multirow{2}{6em}{祭司} &烏利亞的兒子米利末\\
                         & 非尼哈的兒子以利亞撒\\\hline
\multirow{2}{6em}{利未人}&耶書亞的兒子約撒拔\\
                         &賓內的兒子挪亞底 \\\hline
\end{tabular}
\caption{神殿裡接受金銀和器皿的人}
\end{table}

\subsection{亞哈瓦安返耶路撒冷，登寫攜帶的聖殿禮物 8:31-34}
正月十二日，我們從亞哈瓦河邊起行，要往耶路撒冷去。我們神的手保佑我們，救我們脫離仇敵和路上埋伏之人的手。 32 我們到了耶路撒冷，在那裡住了三日。 33 第四日，在我們神的殿裡把金銀和器皿都稱了，交在祭司烏利亞的兒子米利末的手中，同著他有非尼哈的兒子以利亞撒，還有利未人耶書亞的兒子約撒拔和賓內的兒子挪亞底。 34 當時都點了數目，按著分量寫在冊上。

\subsubsection{歸回後向聖殿獻祭 8:31-36}
\begin{table}[H]
\centering 
\begin{tabular}{p{1.5cm}|p{1.5cm}|p{1.5cm}|p{1.5cm}}\hline\hline
公牛 & 12 &公綿羊& 96 \\\hline
綿羊羔& 77&公山羊 &12  \\\hline
\end{tabular}
%\caption{第二次歸回向聖殿獻的祭物}
\end{table}
從擄到之地歸回的人向以色列的神獻燔祭，就是為以色列眾人獻公牛十二隻、公綿羊九十六隻、綿羊羔七十七隻，又獻公山羊十二隻做贖罪祭，這都是向耶和華焚獻的。 36 他們將王的諭旨交給王所派的總督與河西的省長，他們就幫助百姓，又供給神殿裡所需用的。

\section{認罪悔改 9-10}
\subsection{百姓犯與異族通婚的罪 9:1-4}
\begin{table}[!hbp]
\centering 
\begin{tabular}{|p{2cm}|p{2cm}|p{2cm}|p{2cm}|}
\hline
\hline
迦南人&赫人&比利洗人&耶布斯人\\\hline
亞捫人&摩押人&埃及人&亞摩利人\\\hline
\end{tabular}
\caption{外邦女子的種類}
\end{table}
這事做完了，眾首領來見我，說：「以色列民和祭司並利未人沒有離絕迦南人、赫人、比利洗人、耶布斯人、亞捫人、摩押人、埃及人、亞摩利人，仍效法這些國的民行可憎的事。 2 因他們為自己和兒子娶了這些外邦女子為妻，以致聖潔的種類和這些國的民混雜，而且首領和官長在這事上為罪魁。」 3 我一聽見這事，就撕裂衣服和外袍，拔了頭髮和鬍鬚，驚懼憂悶而坐。 4 凡為以色列神言語戰兢的，都因這被擄歸回之人所犯的罪聚集到我這裡來。我就驚懼憂悶而坐，直到獻晚祭的時候。


\subsection{以斯拉的禱告 9:5-15}
獻晚祭的時候我起來，心中愁苦，穿著撕裂的衣袍，雙膝跪下向耶和華我的神舉手， 6 說：「我的神啊，我抱愧蒙羞，不敢向我神仰面，因為我們的罪孽滅頂，我們的罪惡滔天。 7 從我們列祖直到今日，我們的罪惡甚重。因我們的罪孽，我們和君王、祭司都交在外邦列王的手中，殺害、擄掠、搶奪、臉上蒙羞，正如今日的光景。 8 現在耶和華我們的神暫且施恩於我們，給我們留些逃脫的人，使我們安穩如釘子釘在他的聖所，我們的神好光照我們的眼目，使我們在受轄制之中稍微復興。 9 我們是奴僕，然而在受轄制之中，我們的神仍沒有丟棄我們，在波斯王眼前向我們施恩，叫我們復興，能重建我們神的殿，修其毀壞之處，使我們在猶大和耶路撒冷有牆垣。 10 我們的神啊，既是如此，我們還有什麼話可說呢？因為我們已經離棄你的命令， 11 就是你藉你僕人眾先知所吩咐的說：『你們要去得為業之地是汙穢之地，因列國之民的汙穢和可憎的事，叫全地從這邊直到那邊滿了汙穢。 12 所以不可將你們的女兒嫁他們的兒子，也不可為你們的兒子娶他們的女兒，永不可求他們的平安和他們的利益，這樣你們就可以強盛，吃這地的美物，並遺留這地給你們的子孫永遠為業。』 13 神啊，我們因自己的惡行和大罪遭遇了這一切的事，並且你刑罰我們輕於我們罪所當得的，又給我們留下這些人。 14 我們豈可再違背你的命令與這行可憎之事的民結親呢？若這樣行，你豈不向我們發怒，將我們滅絕，以致沒有一個剩下逃脫的人嗎？ 15 耶和華以色列的神啊，因你是公義的，我們這剩下的人才得逃脫，正如今日的光景。看哪，我們在你面前有罪惡，因此無人在你面前站立得住。」

\subsection{以斯拉與眾首領商議聚集以色列人在耶路撒冷聚會 10:1-8}
\subsubsection{在神殿前使祭司長和利未人，以色列眾人起誓 10:1-6}
以斯拉禱告、認罪、哭泣，俯伏在神殿前的時候，有以色列中的男女、孩童聚集到以斯拉那裡，成了大會，眾民無不痛哭。 2 屬以攔的子孫，耶歇的兒子示迦尼對以斯拉說：「我們在此地娶了外邦女子為妻，干犯了我們的神，然而以色列人還有指望。 3 現在當與我們的神立約，休這一切的妻，離絕她們所生的，照著我主和那因神命令戰兢之人所議定的，按律法而行。 4 你起來，這是你當辦的事，我們必幫助你，你當奮勉而行。」

5 以斯拉便起來，使祭司長和利未人並以色列眾人起誓說，必照這話去行。他們就起了誓。 6 以斯拉從神殿前起來，進入以利亞實的兒子約哈難的屋裡，到了那裡不吃飯，也不喝水，因為被擄歸回之人所犯的罪心裡悲傷。

\subsubsection{在約哈難的屋裡定出措施 10:7-8}
他們通告猶大和耶路撒冷被擄歸回的人，叫他們在耶路撒冷聚集。 8 凡不遵首領和長老所議定三日之內不來的，就必抄他的家，使他離開被擄歸回之人的會。

\subsection{會眾同意離棄外邦女子 10:9-17}
\begin{table}[H]
\centering 
\begin{tabular}{p{2cm}|p{8cm}}\hline\hline
九月二十日&眾人都坐在神殿前的寬闊處,會眾同意離棄外邦女子\\\hline
十月初一&一同在座查辦這事\\\hline
正月初一日&查清娶外邦女子的人數\\\hline
\end{tabular}
\caption{查清娶外邦女子的過程}
\end{table}
\subsubsection{以斯拉勸民認罪，離棄異族之妻 10:9-11}
於是，猶大和便雅憫眾人三日之內都聚集在耶路撒冷。那日正是九月二十日，眾人都坐在神殿前的寬闊處；因這事，又因下大雨，就都戰兢。 10 祭司以斯拉站起來，對他們說：「你們有罪了，因你們娶了外邦的女子為妻，增添以色列人的罪惡。 11 現在當向耶和華你們列祖的神認罪，遵行他的旨意，離絕這些國的民和外邦的女子。」

\subsubsection{會眾願遵勸言而行 10:9-14}
會眾都大聲回答說：「我們必照著你的話行。 13 只是百姓眾多，又逢大雨的時令，我們不能站在外頭；這也不是一兩天辦完的事，因我們在這事上犯了大罪。 14 不如為全會眾派首領辦理，凡我們城邑中娶外邦女子為妻的，當按所定的日期，同著本城的長老和士師而來，直到辦完這事，神的烈怒就轉離我們了。」 

\subsubsection{攔阻（總辦）查清娶外邦女子的人員 10:15}
\begin{table}[H]
\centering 
\begin{tabular}{p{10.5cm}|p{6cm}}\hline\hline
 唯有亞撒黑的兒子約拿單、特瓦的兒子雅哈謝阻擋（總辦）這事，&並有米書蘭和利未人沙比太幫助他們。\\\hline
\end{tabular}
%\caption{攔阻查清娶外邦女子的人員}
\end{table}


\subsubsection{查清娶外邦女子的人數 10:16-17}
被擄歸回的人如此而行。祭司以斯拉和些族長按著宗族，都指名見派，在十月初一日，一同在座查辦這事。 17 到正月初一日，才查清娶外邦女子的人數。






\subsection{娶外邦女子的名單 10:18-44}
\subsubsection{祭司中娶外邦女子的名單 10:18-22}
\begin{table}[H]
\centering 
\begin{tabular}{p{3cm}|p{10cm}}\hline\hline
耶書亞的子孫& 約薩達的兒子，和他弟兄瑪西雅、以利以謝、雅立、基大利 \\\hline
音麥的子孫 &哈拿尼、西巴第雅 \\\hline
哈琳的子孫&瑪西雅、以利雅、示瑪雅、耶歇、烏西雅\\\hline
巴施戶珥的子孫&以利約乃、瑪西雅、以實瑪利、拿坦業、約撒拔、以利亞撒\\\hline
\end{tabular}
%\caption{祭司中娶外邦女子的名單}
\end{table}

18 在祭司中查出娶外邦女子為妻的，就是耶書亞的子孫約薩達的兒子和他弟兄瑪西雅、以利以謝、雅立、基大利， 19 他們便應許必休他們的妻。他們因有罪，就獻群中的一隻公綿羊贖罪。 20 音麥的子孫中有哈拿尼、西巴第雅； 21 哈琳的子孫中有瑪西雅、以利雅、示瑪雅、耶歇、烏西雅； 22 巴施戶珥的子孫中有以利約乃、瑪西雅、以實瑪利、拿坦業、約撒拔、以利亞撒。


\subsubsection{利未人中娶外邦女子的名單 10:23-24}
\begin{wraptable}{r}{7.25cm}
\centering 
\begin{tabular}{p{1.9cm}|p{5.cm}}\hline\hline
 利未人 & 約撒拔、示每、基拉雅（基利他），毗他希雅、猶大、以利以謝 \\\hline
歌唱的人  &以利亞實 \\\hline
守門的人  &沙龍、提聯、烏利 \\\hline
\end{tabular}
%\caption{利未人中娶外邦女子的名單}
\end{wraptable}
利未人中，有約撒拔、示每、基拉雅（基拉雅就是基利他），還有毗他希雅、猶大、以利以謝。24 歌唱的人中，有以利亞實。守門的人中，有沙龍、提聯、烏利。


\subsubsection{以色列人中娶外邦女子的名單 10:25-44}
\begin{table}[H]
\centering 
\begin{tabular}{p{2.9cm}|p{12.5cm}}\hline\hline
 巴錄的子孫 & 拉米、耶西雅、瑪基雅、米雅民、以利亞撒、瑪基雅、比拿雅 \\\hline
以攔的子孫  &瑪他尼、撒迦利亞、耶歇、押底、耶利末、以利雅 \\\hline
 薩土的子孫 & 以利約乃、以利亞實、瑪他尼、耶利末、撒拔、亞西撒 \\\hline 
比拜的子孫  &約哈難、哈拿尼雅、薩拜、亞勒 \\\hline
 巴尼的子孫 & 米書蘭、瑪鹿、亞大雅、雅述、示押、耶利末 \\\hline
 巴哈摩押的子孫  &阿底拿、基拉、比拿雅、瑪西雅、瑪他尼、比撒列、賓內、瑪拿西 \\\hline
 哈琳的子孫 & 以利以謝、伊示雅、瑪基雅、示瑪雅、西緬、便雅憫、瑪鹿、示瑪利雅\\\hline
哈順的子孫  &瑪特乃、瑪達他、撒拔、以利法列、耶利買、瑪拿西、示每 \\\hline
 巴尼的子孫 & 瑪玳、暗蘭、烏益、比拿雅、比底雅、基祿、瓦尼雅、米利末、以利亞實、瑪他尼、瑪特乃、雅掃、巴尼、賓內、示每、 示利米雅、拿單、亞大雅、瑪拿底拜、沙賽、沙賴、亞薩利、示利米雅、示瑪利雅、沙龍、亞瑪利雅、約瑟\\\hline
 尼波的子孫  &耶利、瑪他提雅、撒拔、西比拿、雅玳、約珥、比拿雅\\\hline
\end{tabular}
%\caption{利未人中娶外邦女子的名單}
\end{table}
以色列人，巴錄的子孫中有拉米、耶西雅、瑪基雅、米雅民、以利亞撒、瑪基雅、比拿雅； 26 以攔的子孫中有瑪他尼、撒迦利亞、耶歇、押底、耶利末、以利雅； 27 薩土的子孫中有以利約乃、以利亞實、瑪他尼、耶利末、撒拔、亞西撒； 28 比拜的子孫中有約哈難、哈拿尼雅、薩拜、亞勒； 29 巴尼的子孫中有米書蘭、瑪鹿、亞大雅、雅述、示押、耶利末； 30 巴哈摩押的子孫中有阿底拿、基拉、比拿雅、瑪西雅、瑪他尼、比撒列、賓內、瑪拿西； 31 哈琳的子孫中有以利以謝、伊示雅、瑪基雅、示瑪雅、西緬、 32 便雅憫、瑪鹿、示瑪利雅； 33 哈順的子孫中有瑪特乃、瑪達他、撒拔、以利法列、耶利買、瑪拿西、示每； 34 巴尼的子孫中有瑪玳、暗蘭、烏益、 35 比拿雅、比底雅、基祿、 36 瓦尼雅、米利末、以利亞實、 37 瑪他尼、瑪特乃、雅掃、 38 巴尼、賓內、示每、 39 示利米雅、拿單、亞大雅、 40 瑪拿底拜、沙賽、沙賴、 41 亞薩利、示利米雅、示瑪利雅、 42 沙龍、亞瑪利雅、約瑟； 43 尼波的子孫中有耶利、瑪他提雅、撒拔、西比拿、雅玳、約珥、比拿雅。 44 這些人都娶了外邦女子為妻，其中也有生了兒女的。


\section{被掳归回的历史}
%\begin{table}[!hbp]
%\centering 
%\begin{tabular}{|p{2cm}|p{1.5cm}|p{1.2cm}|p{1.5cm}|p{6.cm}|p{2.8cm}|}
%\hline
%\hline
% 外邦人的王 &	做王时间	 &国家	 &参考出处 &	歷史事件 &备注 \\\hline
%那波帕拉沙&	626(630?)-605&	巴比倫&\cite{Nabopolassar}	&	新巴比倫興起 &那波帕拉薩爾 Nabopolassar\\\hline
%	&609		&	&王下23：31-34，代下36：1-4 & 约哈斯被法老擄到埃及& \\\hline
%%	&609	&	&王下23：34,代下36：4&	以利雅敬作猶大和耶路撒冷的王，改名叫約雅敬&約哈斯的哥哥 \\\hline
%尼布甲尼撒&	605-561&	巴比倫&	&	& 那波帕拉沙兒子 Nebuchadnezzar \\\hline
%		&605	&	&耶25：1&	耶利米在约雅敬第四年预言犹大要被掳至巴比伦七十年& \\\hline
%	&	598	&		&	&	约雅敬（609-598）被掳至巴比伦	& \\\hline
%	&586	&&	代下36：11&	西底家（Zedekiah 597-586BC）被掳巴比伦&\\\hline
%以未米罗达& 562-560	& 	&		\\\hline
%\end{tabular}
%\caption{被掳归回的历史}
%\end{table}
 
 \begin{table}[!hbp]
\centering 
\begin{tabular}{|p{3.cm}|p{1.5cm}|p{2.5cm}|p{6.2cm}|p{2.5cm}|}
\hline
\hline
巴比倫王 & 做王时间 &参考出处 & 歷史事件 &备注 \\\hline
那波帕拉沙 Nabopolassar& 626(630?)-605& \cite{Nabopolassar} & 新巴比倫興起,本是一个迦勒底王子 &“那波帕拉薩爾” “拿布普拉撒” \\\hline
&609 &王下23：29，代下35：24 & 約西亞王被埃及王尼哥在米吉多所杀& \\\hline
&609 &王下23：31-34，代下36：1-4 & 约哈斯被法老擄到埃及& \\\hline
&609 & 王下23：34,代下36：4& 以利雅敬作猶大和耶路撒冷的王，改名叫約雅敬&約哈斯的哥哥 \\\hline
&606& 但1：1-2&尼布甲尼撒將耶路撒冷圍困,把神殿中器皿的幾分帶到他神的庫中&\\\hline
&606 & 但1：3-6&但以理,哈拿尼雅、米沙利、亞撒利雅被掳到巴比伦&\\\hline
\textbf{尼布甲尼撒 Nebuchadnezzar}& 605-561&  & & 那波帕拉沙兒子 \\\hline
%&605 & &耶25：1& 耶利米在约雅敬第四年预言犹大要被掳至巴比伦七十年& \\\hline
%&604 & &但2:1& 但以理给尼布甲尼撒解梦& \\\hline
%&?  & &但3:1& 尼布甲尼撒造金像& \\\hline
%& ? & &但4：1& 但以理给尼布甲尼撒解第二个梦& \\\hline
%&  ?& &但4：29, 33& 尼布甲尼撒被趕出離開世人& \\\hline
& 598 & 代下36：6 & 约雅敬（609-598）被掳至巴比伦 & 第一次被掳\\\hline
& 597 &  代下36：10& 尼布甲尼撒差遣人將約雅斤（以西结）和耶和華殿裡寶貴器皿帶到巴比倫 &第二次被掳 \\\hline
& 597 & 代下36：10& 西底家作猶大和耶路撒冷的王 & \\\hline
&586 & 代下36：11& 西底家（Zedekiah 597-586BC）被掳巴比伦&第三次被掳\\\hline
\textbf{以未米罗达 Evilmerodach}& 562-560 && & 尼布甲尼撒之子\\\hline
	&562	&	王下25：27&	把约雅斤（Jehoiachin）释放& \\\hline
%以未米罗达Evil-Merodach && [尼布甲尼撒其子亚米麦督（Amel-Marduk）]	562-560		 & \\\hline	
涅里格利沙爾 Neriglissar & 560-556 & &在前560年殺害了以未米羅達而登位 & 尼布甲尼撒的女婿\cite{Neriglissar} “尼甲沙利薛”\cite{尼甲沙利薛}\\\hline
拉巴施馬爾杜克 Labashi-Marduk&  556& & 巴比倫涅里格利沙爾之子，幼年繼位，在位不足一年便被殺害\cite{Labashi-Marduk}&\\\hline
拿波尼度 Nabonidus&	555-539&	&那波尼德大約在前549年離開巴比倫，任命其子伯沙撒為共同攝政\cite{Nabonidus}& 拿波尼度娶了尼布甲尼撒女儿李道葵斯\cite{尼甲沙利薛}\\\hline
\textbf{伯沙撒	Belshazzar}&553-539&	但5：11，18& & 伯沙撒的母亲相传是尼布甲尼撒的女儿\cite{尼甲沙利薛}\\\hline	

	%		耶51：11	巴比伦帝国要被瑪玳所灭“耶和华定意攻击巴比伦，将她毁灭，所以激动了玛代（米底亚）君王的心，因这是耶和华报仇，就是为自己的殿报仇。”
&553&但7& 伯沙撒元年,但以理見四兽的異象& \\\hline	
	&551\cite{伯沙撒}&但8&伯沙撒第三年,但以理見末后四国異象& \\\hline
		&539  &	 	 但5& 伯沙撒被殺當夜,但以理预言亡国& \\\hline	
\end{tabular}
\caption{被掳归回的历史 巴比倫}
\end{table}

 
\begin{table}[!hbp]
\centering 
\begin{tabular}{|p{2.5cm}|p{1.7cm}|p{1.5cm}|p{6.cm}|p{2.8cm}|}
\hline
\hline
玛代/波斯王 & 做王时间 &参考出处 & 歷史事件 &备注 \\\hline
\textbf{大利烏（玛代）}&	539-538	&	 	 &大利乌成了古列分封的巴比伦王\cite{伯沙撒}	&   \\\hline
&	元年538	&	但9	 &   但以理得知耶路撒冷荒涼七十年。便禁食，向神祈禱得知七十個七的启示&   \\\hline
&	元年538	&	但11：1	 &   我曾起來扶助米迦勒，使他堅強。&   \\\hline
 &	 	&	但6：28	 &   但以理当大利乌王在位的时候和波斯王塞鲁士在位的时候，大享亨通。	&   \\\hline
&	 	&	但6	 &   但以理被投入狮子坑	&   \\\hline
\textbf{古列 Cyrus    （波斯）}&	559-529	&代下36: 22 拉1:1 && 又名“塞鲁士”， “居魯士二世” \\\hline	
&538（7）\cite{居鲁士}&拉1:1 & 波斯王古列元年下诏，重建圣殿&   \\\hline	
&536&但10-12 & 塞鲁士第三年，但以理得末后大争战的启示。&   \\\hline	
&536& &所罗巴伯回归耶路撒冷&   \\\hline	
冈比西斯二世 Cambyses&	529-522	&  &居鲁士大帝的儿子，冈比西斯二世将兄弟巴尔迪亚杀死继王位  &   \\\hline	
高墨达\cite{高墨达}&	522（3个月）&& 伪称自己是巴尔迪亚，夺取了王位 &斯梅尔迪士，  亚达薛西    \\\hline	
\textbf{大流士大帝 Darius I，（Darius the Great）}& 522-486&  &居鲁士大帝的女婿&  又名“大利乌”，“大流士一世” \\\hline	
& 520&&大利乌王下旨继续建殿的工作& \\\hline
&516&拉6：15&殿在大流士六年亚达月初三完成& \\\hline
薛西一世  & 486 &&&Xerxes I \\\hline
\textbf{亞哈隨魯}&	485-465	&	 \cite{亞哈隨魯}拉四：6 &亚哈随鲁登基，敌人上本控告犹大和耶路撒冷的居民。大流士一世与居鲁士大帝之女阿托莎的儿子& “薛西斯一世” “澤克西斯一世” “泽尔士一世” \\\hline	
&479(8)　& & 以斯帖成为王后&\\\hline
\textbf{亚达薛西一世” Artaxerxes I} &	465-425&拉四：7-16&澤克西斯一世之子&阿尔塔薛西斯一世\\\hline
&458&&以斯拉回归耶路撒冷	&\\\hline
&444&&尼希米到达耶路撒冷，城墙完工	&\\\hline
薛西斯二世	&425&&阿尔塔薛西斯一世的嫡子，在位45天被暗杀&\\\hline	
塞基狄亚努斯 	&425&&据说他是阿尔塔薛西斯一世的庶子，仅仅在位6个月零15天便被大流士二世杀死&\\\hline	
大流士二世	&423-404&&夺取异母兄弟塞基狄亚努斯王位&\\\hline
&420&&尼希米再到耶路撒冷	&\\\hline
%拉六：3，14  塞鲁士王元年，他降旨论到耶路撒冷上帝的殿，要建造这殿为献祭之处，坚立殿的根基。殿高六十肘，宽六十肘。。。犹大长老因先知哈该和易多的孙子撒迦利亚所说劝勉的话，就建造这殿，凡事亨通。他们遵着以色列上帝的命令和波斯王塞鲁士、大利乌（或大流士一世 Darius I or Darius the Great，公元前522-前486年)、亚达薛西（Artaxerxes I，公元前465-前425年)的旨意，建造完毕。								
%在先知书上，我们也可以看到他的名字：									
%但一：21  到塞鲁士王元年，但以理还在。									
								
%但十：1  波斯王塞鲁士第三年，有事显给称为伯提沙撒的但以理。这事是真的，是指着大争战。但以理通达这事，明白这异象。														
%但最令人惊讶的是，先知以赛亚竟然预言塞鲁士在历史舞台上的出现，这是不可%思议的事，因为以赛亚是在“乌西雅（Uzziah 776/775-736/735BC）、约坦（Jotham 750-735/730BC）、亚哈斯（Ahaz 735/734 or 731/730-715BC）、希西家（Hezekiah 715-687/686BC）作犹大王的时候得默示”（赛一：1）这是在塞鲁士（Cyrus II 559-529BC）出现之前约两百年。有可能吗？不信派的人当然说不可能，特别是名字都出现在《以赛亚书》三十九章之后：									
%赛四十四：28  论塞鲁士说：“他是我的牧人，必成就我所喜悦的，必下令建造耶路撒冷，发命立稳圣殿的根基。”									
%赛四十五：1  我耶和华所膏的塞鲁士，我搀扶他的右手，使列国降伏在他面前。我也要放松列王的腰带，使城门在他面前敞开，不得关闭。我对他如此说：。。						赛四十五：13  我凭公义兴起塞鲁士，又要修直他一切道路。他必建造我的城，释放我被掳的民，不是为工价，也不是为赏赐。这是万军之耶和华说的。						
%还有赛四十一：2-3，25-26，四十六：11 都暗示着“塞鲁士”的到来。								
\end{tabular}
\caption{被掳归回的历史2}
\end{table}


\begin{thebibliography}{1}
  \bibitem{Nabopolassar} https://zh.wikipedia.org/wiki/%E9%82%A3%E6%B3%A2%E5%B8%95%E6%8B%89%E8%96%A9%E7%88%BE
  \bibitem{Neriglissar} https://zh.wikipedia.org/wiki/%E6%B6%85%E9%87%8C%E6%A0%BC%E5%88%A9%E6%B2%99%E7%88%BE

  \bibitem{Labashi-Marduk} https://zh.wikipedia.org/wiki/%E6%8B%89%E5%B7%B4%E6%96%BD%E9%A6%AC%E7%88%BE%E6%9D%9C%E5%85%8B
  
  \bibitem{Nabonidus}  https://zh.wikipedia.org/wiki/%E9%82%A3%E6%B3%A2%E5%B0%BC%E5%BE%B7
  
  \bibitem{亞哈隨魯}   https://zh.wikipedia.org/wiki/%E4%BA%9E%E5%93%88%E9%9A%A8%E9%AD%AF

  \bibitem{伯沙撒}  http://www.sdacn.org/htmlshuji/yuyan/daniel-m/9.htm

   \bibitem{尼甲沙利薛}  http://www.jiduzhijia.com/cj/Old%20Testament/27Dan/27TT00.htm
   
   \bibitem{居鲁士}  https://wapbaike.baidu.com/item/%E4%BB%A5%E6%96%AF%E6%8B%89

   \bibitem{高墨达}  https://zh.wikipedia.org/wiki/%E9%AB%98%E5%A2%A8%E8%BE%BE

  \bibitem{notes} John W. Dower {\em Readings compiled for History
  21.479.}  1991.

  \bibitem{impj}  The Japan Reader {\em Imperial Japan 1800-1945} 1973:
  Random House, N.Y.

  \bibitem{norman} E. H. Norman {\em Japan's emergence as a modern
  state} 1940: International Secretariat, Institute of Pacific
  Relations.

  \bibitem{fo} Bob Tadashi Wakabayashi {\em Anti-Foreignism and Western
  Learning in Early-Modern Japan} 1986: Harvard University Press.


  \end{thebibliography}


